%% kapa: exam=1 round=36 year=2025 period=1 form=A subject=civil_law subject_ko=민법 seq=3 number=3
%% source: Q-Net 감정평가사 기출문제 (cst00309) artlSeq=5231525 file=raw/1008/5231525/2025년 제36회 감정평가사 1차시험 1교시 문제지 원본.pdf
%% conversion: kordoc 4.24.51 -> tools/gen1.py (auto); review=proofread-sonnet; flags=none
\begin{question}{3}
제한능력자의 상대방보호에 관한 설명으로 옳지 않은 것은? (다툼이 있으면 판례에 따름)
\begin{choices}
  \item 상대방은 제한능력자의 법정대리인에게 1개월 이상의 기간을 정하여 취소할 수 있는 행위를 추인할 것인지 여부의 확답을 촉구할 수 있다.
  \item 계약 당시에 당사자가 제한능력자임을 알지 못한 상대방은 추인이 있을 때까지 제한능력을 이유로 취소할 수 있는 계약의 의사표시를 철회할 수 있다.
  \item 제한능력자의 취소할 수 있는 단독행위는 추인이 있을 때까지 상대방이 거절할 수 있다.
  \item 상대방은 제한능력을 이유로 취소할 수 있는 계약에 대한 철회를 제한능력자에게 할 수 없다.
  \item 적극적인 속임수를 사용하여 자기를 능력자로 믿게 하여 상대방과 계약을 체결한 제한능력자는 제한능력을 이유로 그 계약을 취소할 수 없다.
\end{choices}
\end{question}
\answer{4}
